%% =====================================================================
%%  T-23-03  Testing a Claimed Credential
%%  -------------------------------------------------------------------
%%  One idea per file. Core is mandatory. Slots in canonical order.
%%  Max 700 words. No layout commands. No filler. New source material
%%  for this entry goes in sources/<BOOK>/T-23-03.tex -- never by
%%  rewriting this file (docs/RULES.md R0.1).
%% =====================================================================
%% @kind: topic
%% @id: T-23-03
%% @title: Testing a Claimed Credential
%% @subtitle: Verification that does not depend on the person being verified
%% @chapter: c23
%% @order: 30
%% @status: stable
%% @sources: B4
%% @updated: 2026-09-19
%% @allow-rewrite: slot migration to the seven-question format (docs/RULES.md section 2)

\topic{T-23-03}{Testing a Claimed Credential}{Verification that does not depend on the person being verified}


\begin{core}
A credential is worth exactly as much as your ability to check it without the claimant's help. Any verification that runs through the person claiming authority --- their number, their reference, their colleague --- verifies nothing at all.
\end{core}

\begin{mechanism}
The failure mode is circular. A claimed authority offers the means of checking the claim, the means are controlled by the claimant, and the check confirms whatever was asserted. The circle is invisible in the moment because the check feels like diligence.

Breaking it requires finding a route to the institution that does not pass through the individual: a number looked up independently, a registry searched by name, a colleague contacted through a switchboard rather than a mobile. Where no such route exists, the honest conclusion is that the credential is unverifiable --- which is a usable finding, not a failure. Unverifiable authority should be treated as absent rather than as presumed.

A second test is available and costs nothing: ask for the claimant's reasoning rather than their conclusion. Genuine expertise can show its work, including what would change it. A symbol repeats the position, becomes irritated at the question, or refers back to the credential itself.
\end{mechanism}

\begin{application}
  \item Establish a route to the institution before making contact, so that the check does not have to be negotiated later.
  \item Ask for reasoning rather than conclusions; it is cheap and it separates the two cases quickly.
  \item Ask what would change their mind, and note whether the answer is about evidence or about loyalty.
  \item Where the credential is unverifiable, say so plainly and proceed as though it were absent.
\end{application}

\begin{feedback}
  \item You are being asked to verify through channels the claimant supplied.
  \item The credential is asserted more often than it is evidenced.
  \item Questions about reasoning are answered with questions about your qualifications to ask.
  \item The claim is time-sensitive, which removes the opportunity to check.
  \item The title is specific but the conferring body is vague.
\end{feedback}

\begin{failure}
Verification is slow, occasionally rude, and frequently unnecessary --- most claimed authority is genuine, and checking all of it makes ordinary life unworkable. The test is therefore targeted: run it where money, consent, medical treatment, legal position or access is at stake, where the contact is unusual in arriving, or where urgency is asserted. Elsewhere accept the symbol and move on.
\end{failure}

\begin{countermeasures}
  \item Look the institution up yourself. Do not use any number, address or contact given to you.
  \item Call back through a switchboard and ask for the person by name. A real employee is reachable that way; a costume is not.
  \item Search a public registry. Most regulated professions maintain one and most claimants assume you will not.
  \item Delay the decision by a day. Urgency is the one condition under which verification becomes impractical, and it is usually manufactured for that reason.
  \item Ask the question in writing. A claim that survives speech often does not survive being written down and dated.
\end{countermeasures}

\sources{B4}
\seesources
\seealso{T-18-01, T-18-02, T-23-01, T-24-03}
